% =====================================================================
%  Syntaktická analýza
% =====================================================================
\section{Syntaktická analýza}\label{sec:parser}

\subsection{Rekurzivní sestup}\label{sec:rek-sestup}
\TODO{Odpovídá parser gramatice funkcí na neterminál? Kolik tokenů čte dopředu?}

\subsection{Precedenční analýza výrazů}\label{sec:precedence}
\TODO{Kdy parser předává řízení precedenční analýze a jak se vrací?}

% --- precedenční tabulka ---------------------------------------------
% \shift <  \reduce >  \equal =  prázdná buňka = chyba
\begin{table}[H]
  \centering
  \renewcommand{\arraystretch}{1.2}
  \begin{tabular}{c|cccc}
           & \TODOinline{op} & \texttt{(} & \texttt{)} & \$ \\ \hline
    \TODOinline{op} & & & & \\
    \texttt{(}      & & & & \\
    \texttt{)}      & & & & \\
    \$              & & & & \\
  \end{tabular}
  \caption{Precedenční tabulka (řádek = terminál na zásobníku,
           sloupec = vstupní token)}
  \label{tab:precedence}
\end{table}

\subsection{Sémantická analýza}\label{sec:semantika}
Pro provedení sémantické analýzy se volá funkce \texttt{semantic\_analyze}.
\begin{itemize}[nosep]
  \item Pass 0: Přidá všechny vestavěné funkce do tabulky symbolů.
  \item Pass 1: Projde všechny funkce a přidá je do tabulky symbolů.
  \item Pass 2: Projde tělo každé funkce a provede sémantickou analýzu,
        včetně kontroly pravidel, typových kontrol, typových konverzí,
        kompatibility typů a odvozování typů. Zároveň do AST doplní
        informace potřebné pro codegen.
\end{itemize}

\TODO{Něco doplnit k průchodům, třeba val a var nebo návratové typy?}

\subsection{Generování cílového kódu}\label{sec:codegen}
\TODO{Jak se řeší rámce, návěští a vestavěné funkce?}
